% ============================================================
% Tabelle 7.3-3: Lückentypen beider Konfigurationen — v01 (08.09.)
% In Overleaf ablegen als: tables/eval-lueckentypen.tex
% Aufruf im Fließtext: \input{\tabledir/eval-lueckentypen.tex}
% Label: t:eval-lueckentypen
% Stil: chap6.1-Konvention.
% HERKUNFT: runs/w2-nachreview-konsolidierung.json §lueckentypen
% (konsolidierter Stand; Typisierung je ID aus den beiden Reviews
% übernommen, drei neue Lücken dokumentiert nachtypisiert:
% Ledger-062=F, 107=O beidseitig, F-089=U). Typ beschreibt den
% FEHLENDEN Bestandteil der jeweiligen Ausgabe — dieselbe Gold-ID
% kann je Konfiguration verschieden typisiert sein.
% ============================================================

\begin{table}[htbp]
  \centering
  \small
  \begin{tabular}{>{\raggedright\arraybackslash}p{0.52\textwidth}rr}
    \hline
    \textbf{Fehlender Bestandteil} & \textbf{Ledger} & \textbf{F} \\
    & (55 Lücken) & (40 Lücken) \\
    \hline
    Funktion / Berechtigung / fachlicher Ablauf & 18 & 6 \\
    Darstellung / konkrete Bediengestaltung & 18 & 18 \\
    Zweck / Qualitätsziel & 11 & 8 \\
    Offene Frage / technische Konkretisierung / Interpretationsfall & 8 & 8 \\
    \hline
    Summe & 55 & 40 \\
    \hline
  \end{tabular}
  \caption[Typisierung der Abdeckungslücken im Stressfall]{Retrospektive Typisierung der unvollständig abgedeckten
    Referenzaussagen des Stressfalls nach dem jeweils fehlenden
    Bestandteil (konsolidierter Auswertungsstand). Der Typ
    beschreibt, \emph{was} in der jeweiligen Ausgabe fehlt --- er
    ist keine Wichtigkeits- oder Schweregradskala, und dieselbe
    Referenzaussage kann je Konfiguration unterschiedlich typisiert
    sein, wenn die Ausgaben unterschiedliche Teile erhalten.}
  \label{t:eval-lueckentypen}
\end{table}
